% !TEX root = ../../../main.tex
% Block card: catalogue.window_images -- "Encode windows" (Part I, the blocks chapter).
\cardfield{What it does}
Makes one picture per window of a window set, the way the pictures for training the image
networks were made, and stacks them ready to be scored.

\cardfield{When to use it}
Between a windows block (such as \blockref{preprocessing.sliding_windows}) and
\blockref{catalogue.cnn_score}, when you want a trained network to say which windows look
interesting. The page shows the pictures exactly as the network will see them.

\cardfield{What goes in, what comes out}
In goes a set of windows on one channel. Out comes a stack of square colour pictures, one per
window, all the same size (224 pixels a side by default).

\cardfield{Settings that matter}
\emph{image\_type} picks the picture: the colour fusion (\blockref{catalogue.gramian_fusion})
or one of the three single images, shown in grey. Pick the one the network you will use was
trained on; the fusion networks need \emph{fusion}. \emph{img\_size} should match the
network's input, 224. \emph{max\_windows} refuses a set so large the stack would not fit in
memory.

\cardfield{Watch out for}
Each window is first drawn at one pixel per sample and only then shrunk, so a long window is
slow and memory-hungry even though the final picture is small: an hour at 10\,Hz is a
picture of 36{,}000 by 36{,}000 before shrinking, and the block does not stop you. The
networks were built for 10-minute windows at 1\,Hz; other lengths are outside what they
saw. A window with a missing sample comes out black.
